% ============================================================================
%  Appendice D -- I sorgenti completi di Verm1
% ============================================================================
\chapter{Verm1: i sorgenti completi}
\label{app:verm1}

Questa appendice stampa per intero i quattro programmi di riferimento del
progetto e i file del bundle dello Store. L'implementazione di riferimento di
\vermone è un contributo di Giuseppe Coviello; i capitoli tutorial sono
costruiti su di essa. Sono i file che i test automatici del materiale del
corso compilano, eseguono sull'esecutore senza interfaccia di \prgqt e giocano
con un giocatore programmato; i listati dei capitoli tutorial sono regioni di
questi stessi file. Se il vostro programma si rifiuta di funzionare all'inizio
di un tutorial, ripartite dal file della parte precedente
(Tabella~\ref{tab:verm1:files}).

\begin{table}[htb]
\centering\small
\caption{I programmi di riferimento.}\label{tab:verm1:files}
\begin{tabularx}{\textwidth}{@{}llrL@{}}
\toprule
\textbf{File} & \textbf{Capitolo} & \textbf{Righe} & \textbf{Il gioco a questo punto}\\
\midrule
\fname{verm1\_part1.S} & \ref{ch:verm1p1} & 78 & una griglia, una testa che si muove da sola, un mondo senza bordi\\
\fname{verm1\_part2.S} & \ref{ch:verm1p2} & 231 & la guida, la regola dell'inversione a U, il cibo, un punteggio a tacche, un suono\\
\fname{verm1\_part3.S} & \ref{ch:verm1p3} & 435 & un corpo che cresce, le collisioni, tre stati, un punteggio migliore\\
\fname{verm1\_part4.S} & \ref{ch:verm1p4} & 565 & cifre, livelli, due suoni, disegno incrementale: il gioco finito\\
\bottomrule
\end{tabularx}
\end{table}

Ogni parte si compila e si esegue con gli stessi due comandi:

\begin{shell}
python3 -m prg32 cartridge build verm1_part4.S --portable \
    --entry-prefix verm1 --name verm1 --out verm1.prg32
python3 -m prg32 esp32c6 upload-and-run verm1.prg32 \
    --url http://localhost:8080 --slot cart0
\end{shell}

\section{Parte 1: una testa che attraversa il mondo}
\asmfile[caption={\fname{verm1\_part1.S}.},label=lst:verm1:p1]{code/verm1/verm1_part1.S}

\section{Parte 2: la guida, il cibo e un punteggio}
\asmfile[caption={\fname{verm1\_part2.S}.},label=lst:verm1:p2]{code/verm1/verm1_part2.S}

\section{Parte 3: un corpo, regole e stati}
\asmfile[caption={\fname{verm1\_part3.S}.},label=lst:verm1:p3]{code/verm1/verm1_part3.S}

\section{Parte 4: il gioco finito}
\asmfile[caption={\fname{verm1\_part4.S}.},label=lst:verm1:p4]{code/verm1/verm1_part4.S}

\section{I file del bundle dello Store}

I tre file che seguono sono quelli del Capitolo~\ref{ch:store}, per il gioco
di riferimento. In un vostro bundle, l'identificatore, il titolo e ogni
occorrenza di \texttt{Your Name} vengono sostituiti; il ringraziamento per
l'implementazione di riferimento resta, come per qualsiasi lavoro che parte da
quello di qualcun altro.

\jsonfile[caption={\fname{manifest.json}.},label=lst:verm1:manifest]{code/verm1/store/manifest.json}

\jsonfile[caption={\fname{metadata.json}: il manifesto senza \code{assets} e \code{architectures}.},label=lst:verm1:metadata]{code/verm1/store/metadata.json}

\jsonfile[caption={\fname{colophon.json}.},label=lst:verm1:colophon]{code/verm1/store/colophon.json}
